\section{A growing branching core with scalar flux events}
\label{sec:op3-branch-core-flux}

\paragraph{Status and purpose.}
This section is a \statusproved{} proof draft awaiting independent review.
It combines the original-residual publication bands of
\cref{sec:op3-geometric-events} with low-degree path elimination.
The number of retained vertices may now grow. A fully specified dense
inverse pays for their response changes, leaving an explicit quadratic
core-size factor. No high-dimensional geometric reporter is required.
General OP3 remains \statusopen.

Set $\lambda=\epsappr/2$, $\bm M=\bm D-\gamma\bm A$ and
$\bm b=\bm e_v-\lambda\bm d$. At each stage, $U$ is the admitted set
and $\bm u^U_U=\bm M_{UU}^{-1}\bm b_U>0$, with zero outside $U$.
Retain the seed and every admitted vertex of original degree at least three
in a set $C$. Immediately eliminate all other admitted vertices into saved
scalar equations. Let $r=|C|$ at termination, $n_U=|U|$, and $V=\vol(U)$.
These are measured properties of the returned support, not input advice.

\subsection{One retained endpoint per unfinished path}

\begin{lemma}[Stable physical-flux arms]
\label{lem:op3-scalar-flux-arms}
For every original edge $ij$ with $i\in U$ and $j\notin U$, the physical
flux $F_{ij}=\gamma u_i^U$ has a representation
\begin{equation}
 F_{ij}=a_{ij}u_{c(ij)}^U+b_{ij},\qquad
 a_{ij}>0,\quad b_{ij}\leq0,\quad c(ij)\in C.
 \label{eq:op3-physical-flux-arm}
\end{equation}
The coefficients and core label are fixed during this edge's lifetime on
the boundary. An admission creates at most one such record per newly scanned
original edge, and retires the records ending at its vertex. The total number
of arm creations and retirements is $O(V)$.
\end{lemma}
\begin{proof}
An eliminated vertex has original degree at most two. Every connected
component of eliminated vertices is consequently a path, and it touches the
retained set: the admitted set is connected to the retained seed. If a path
also touches an inactive vertex, it has exactly one retained endpoint.
It cannot have an additional inactive or retained exit. Components joining
two retained endpoints have no inactive exit and produce no frontier arm.

A direct retained-to-inactive edge has $(a,b)=(\gamma,0)$.
Suppose an admitted degree-two vertex $i$ has an inactive successor $j$.
It has exactly one original active neighbor, so its old frontier record
contains exactly one arm $a u_c+b$. With current Schur diagonal $\delta_i$,
its equation, keeping $j$ as a formal outside variable, is
\[
 \delta_i u_i=-\lambda d_i+b+a u_c+\gamma u_j.
\]
Eliminating $i$ creates the new physical-flux arm, evaluated at $u_j=0$,
\[
 F_{ij}=\frac{\gamma a}{\delta_i}u_c
       +\frac{\gamma(-\lambda d_i+b)}{\delta_i}.
\]
Its diagonal contribution at $j$ is $-\gamma^2/\delta_i$.
Positive Schur pivots preserve the coefficient signs. Nothing changes an
existing arm's internal path until its target is admitted, when that arm
is retired. Other admissions change only its retained endpoint potential.
Each original edge is first seen when one endpoint's row is scanned;
an edge dictionary avoids counting its reverse appearance again.
\end{proof}

Several arms may end at the same inactive vertex, including arms with
different retained endpoints. Their contributions are summed in one
candidate record. In particular, the exact gate is
\begin{equation}
 g_j=-\lambda d_j+\sum_{ij:\,i\in U}F_{ij}.
 \label{eq:op3-arm-original-gate}
\end{equation}
It is not a collection of independent-copy gates. A candidate's stored
Schur diagonal includes all prior path eliminations into that candidate.

\subsection{Scalar queues certify the original residual scale}

For each arm keep a lower published flux $\ell_{ij}$, initially zero.
Use $h=\epsappr/16$ and $\kappa=5/4$. The next publication is due when
\begin{equation}
 u_{c(ij)}^U>
 \tau_{ij}:=\frac{\kappa\ell_{ij}+h-b_{ij}}{a_{ij}}.
 \label{eq:op3-arm-next-event}
\end{equation}
A minimum heap for each retained coordinate stores these thresholds.
After updating the retained potentials, visit each heap and process all
due entries. Publish the actual flux $\ell_{ij}\leftarrow F_{ij}$, update
the candidate's sum $L_j=\sum_i\ell_{ij}$, and insert the new threshold.
Retired arms may be removed lazily: every discarded entry is charged once.

When all heaps are settled, every live arm satisfies
\[
 \ell_{ij}\leq F_{ij}\leq\tfrac54\ell_{ij}+\epsappr/16.
\]
Queue a candidate as soon as $L_j>(11/20)\epsappr d_j$. This is a legal
admission because its exact gate in \cref{eq:op3-arm-original-gate} is
strictly positive for $\lambda=\epsappr/2$. A queued candidate stays legal:
all restricted solutions and physical fluxes increase under legal admissions.
When the ready queue is empty after settling all heaps,
\begin{equation}
 0\leq r_j^U=\sum_i F_{ij}
 \leq\tfrac54L_j+\tfrac1{16}\epsappr d_j
 \leq\tfrac34\epsappr d_j
 \quad (j\notin U).
 \label{eq:op3-arm-terminal-residual}
\end{equation}
An unexposed vertex has no active neighbor and residual zero.
Active residuals are exactly $\lambda d_i$. Thus the exact restricted
solution already supplies an ACL output; a further solve is unnecessary
in this exact-arithmetic implementation.

\begin{lemma}[All scalar reporting work is paid]
\label{lem:op3-scalar-flux-work}
Write $L=1+\log_+(1/(\bar\alpha\epsappr))$. Over the entire run the arm
queues perform $O(VL)$ publications, heap insertions and heap removals,
plus $O(n_Ur+VL)$ top inspections. Their work, including candidate sums,
ready entries and failed event searches, is
$O(n_Ur+VL\log(2+V))$, with $O(V+r)$ stored words.
\end{lemma}
\begin{proof}
Conservation on any positive restricted face gives
$\bar\alpha\sum_{i\in U}d_i u_i^U\leq1$, hence
$F_{ij}\leq\gamma/\bar\alpha$. A first publication exceeds $h$;
each later one exceeds $5/4$ times its predecessor. Each fixed arm therefore
publishes $O(L)$ times. There are $O(V)$ arms by
\cref{lem:op3-scalar-flux-arms}. A publication pops its current entry and
pushes exactly one replacement. Retiring an arm leaves at most one stale
entry. Thus all stale removals and allocated heap records are paid.
Each round scans at most $r$ heaps, and every round except the last admits
one vertex. One failed top comparison per nonempty heap is included.
Every candidate enters the ready queue at most once. All arm records,
candidate records and possibly stale heap entries occupy $O(V+r)$ words.
\end{proof}

This argument uses individual physical fluxes, not relative error after
division by a candidate's accumulated negative Schur load. The latter
contract fails on the source-valid path examples in
\texttt{SCHUR\_RELATIVE\_REPORTER\_PROBE.md}. Large cancellation in an
arm's affine formula is allowed in the exact-real word model; coefficient
bit cost and finite-precision certification remain outside this theorem.

\subsection{An explicit core backend and the resulting bound}

Maintain the inverse $\bm P=\bm H_C^{-1}$ and retained values $\bm u_C$.
At an admission the candidate has a Schur diagonal $\delta$, load $b<0$
and nonnegative sparse coupling vector $\bm a$. Put
\begin{equation}
 \bm z=\bm P\bm a,\qquad s=\delta-\bm a^{\mathsf T}\bm z>0,
 \qquad t=(b+\bm a^{\mathsf T}\bm u_C)/s>0.
 \label{eq:op3-core-rank-update}
\end{equation}
All old retained values change by $\bm u_C\leftarrow\bm u_C+t\bm z$.
If the candidate has degree at most two, eliminate it and update
$\bm P\leftarrow\bm P+\bm z\bm z^{\mathsf T}/s$.
If it is a branching vertex, retain it with value $t$ and bordered inverse
\[
 \begin{pmatrix}
 \bm P+\bm z\bm z^{\mathsf T}/s&\bm z/s\\
 \bm z^{\mathsf T}/s&1/s
 \end{pmatrix}.
\]
These are ordinary scalar block-inverse identities. A low-degree candidate
has at most two arm coordinates; at a branching admission, forming its
couplings costs at most its original degree. Every dense inverse entry
updated, every retained value written and every inverse allocation is
charged. A reverse pass through the saved low-degree equations recovers
all eliminated values in $O(V)$ work.

\begin{theorem}[Local ACL with a growing branching-core parameter]
\label{thm:op3-branch-core-flux}
On every graph in the canonical local access model, the described algorithm
returns a nonnegative ACL approximation at tolerance $\epsappr$. It scans
precisely the rows of its positive output set $U$ and queries only the seed
degree and degrees of exposed neighbors. For a nonempty output, its exact-real
word work is bounded by
\begin{equation}
 O\left(n_Ur^2+rV+VL\log(2+V)\right),\qquad
 L=1+\log_+\frac1{\bar\alpha\epsappr},
 \label{eq:op3-branch-core-work}
\end{equation}
and its working storage is $O(V+r^2)$. Empty output costs $O(1)$.
Moreover $V\leq2/\epsappr$, so the work is
$\widetilde O((1+r^2)/\epsappr)$. In particular, this attains the OP3
soft-$O(1/\epsappr)$ scale when the returned support has at most a
polylogarithmic number of original branching vertices.
\end{theorem}
\begin{proof}
If $\lambda d_v\geq1$, zero is feasible with the required residual bound.
Otherwise the positive singleton seed initializes the construction. The
Schur identities and \cref{lem:op3-scalar-flux-arms} preserve exact restricted
solutions. Strict positive gates give positive new coordinates and monotone
old ones by inverse positivity. As in \cref{thm:op3-two-port-frontier}, every
admission belongs to the positive support of the full obstacle at $\lambda$.
Thus the process terminates and the finite-band gate rule is safe.
Conservation gives $\lambda V\leq1$ on every reached face.

At termination, \cref{eq:op3-arm-terminal-residual} and active-face
equalities give $0\leq\bm e_v-\bm M\bm u^U\leq\epsappr\bm d$.
Consequently $\widehat{\bm\pi}=\bar\alpha\bm D\bm u^U$ has the required
ACL witness. The final reverse pass evaluates every eliminated equation
once and writes only positive output coordinates.

Each of $n_U-1$ admissions changes at most $r^2$ inverse entries. Sparse
inverse-vector products cost at most $O(rV)$ altogether; the sum of
candidate arm counts is $O(V)$. Dense inverse growth allocates $O(r^2)$
entries. Original row entries, cached degrees, edge dictionaries and stored
reverse records cost $O(V\log(2+V))$ using comparison dictionaries.
\Cref{lem:op3-scalar-flux-work} pays for every event and failed queue test.
These bounds include initialization and combine to
\cref{eq:op3-branch-core-work}. Only the inverse is quadratic in storage.
\end{proof}

\paragraph{Implementation and remaining lemma.}
\texttt{branch\_core\_flux.py} implements this construction with exact
fractions, scalar binary heaps and explicitly updated dense inverse entries.
Its Python dictionaries are an implementation convenience; comparison
dictionaries give the conservative word bound above. Full-graph solves,
inverse checks and residual checks occur only in the separate audit wrapper.
This implementation does not reproduce the full obstacle support: its output
is an exact restricted solution chosen to meet the ACL residual tolerance.

The remaining general-graph obligation is now concentrated in the retained
core: replace dense inverse changes and full scans of the growing core's
queues by an implicit response primitive that reports all due scalar
thresholds and certifies quietness. A fast named coordinate query or a fresh
near-linear core solve alone does not pay for that sequence. The polynomial
$r^2$ factor is an upper bound for this backend, not a lower bound for OP3.

\begin{proposition}[Dense-core cost on an easy tree family]
\label{prop:op3-dense-core-binary-tree}
On the complete binary tree of height $h\geq2$, take seed at the root,
$\alpha=1/1009$ and $\epsappr=\alpha/(10\cdot6^h)$.
Every ACL output has full support. The explicit inverse backend above
therefore performs $\Omega(n^3)$ word writes on this $n=2^{h+1}-1$ vertex
family, while its total number of flux publications is near-linear up to
logarithmic factors. Scanning all retained queues is separately charged.
This is a statement about the explicit retained-inverse representation.
\end{proposition}
\begin{proof}
The Neumann expansion of $\bm M^{-1}$ contains the unique source-to-vertex
walk of length at most $h$. Since all degrees are at most three and
$\gamma>1/2$, its contribution gives
\[
 \frac{\pi_i}{d_i}
 =\bar\alpha(\bm M^{-1})_{iv}
 \geq\frac{\bar\alpha}{3\cdot6^h}
 >\epsappr.
\]
ACL implies degree-normalized semantic error at most $\epsappr$, so no
coordinate can be zero in an ACL output. The retained set therefore has
$r=2^h-1$ vertices. At each successive retained birth, the previous inverse
has order $k=1,\ldots,r-1$; its top-left update changes every one of its
$k^2$ entries. Indeed the current restricted graph is connected, its
retained Schur inverse is entrywise positive, and the new coupling is
nonzero and nonnegative, so the update vector $\bm z$ is strictly positive.
The mandatory writes total
$\sum_{k=1}^{r-1}k^2=r(r-1)(2r-1)/6=\Omega(n^3)$.
The same reasoning forces changes to all explicitly stored old retained
values, separately from inverse maintenance. Balanced trees already have
short rooted paths; this is not a graph-access hardness example.
\end{proof}
